% Source: MQ15b_v1.html
\documentclass[11pt]{article}
\usepackage[margin=1in]{geometry}
\usepackage{amsmath}
\usepackage{amssymb}
\usepackage{siunitx}
\usepackage{enumitem}

\title{MQ15b: Waves}
\author{}
\date{}

\begin{document}
\maketitle

\section*{MQ15b: Waves}

\begin{enumerate}[leftmargin=*,label=\textbf{Question \arabic*.}]
\item \hfill \textit{1 pts}

Two traveling waves of the same frequency (\emph{f} = 200 Hz) and wavelength interfere to produce a standing wave. The distance from node \#3 to node \#5 is 5.27 m. What is the wavelength (in m) of the original waves?

\item \hfill \textit{1 pts}

A vibrating string has a fifth harmonic of 1,436 Hz. What is the second harmonic in Hertz?

\item \hfill \textit{1 pts}

A wire with mass 0.04 kg and length 0.57 m has a fundamental frequency of 105 Hz. Calculate the tension in the wire in Newtons.

\item \hfill \textit{1 pts}

A wire with mass 0.07 kg and length 0.67 m is under tension of 516 N. Calculate the wire's fundamental frequency in Hz.

\item \hfill \textit{1 pts}

A given wave will have a certain number of wavelengths per meter. The number of radians corresponding to that number of wavelengths is called the [blank].

\item \hfill \textit{1 pts}

A wave whereby the vibration is perpendicular to the direction of propagation is called a [Select]
 wave, whereas a wave having the vibration parallel to the direction of motion is called a [Select]
 wave.

\item \hfill \textit{1 pts}

A standing wave on a vibrating string has points of totally destructive interference called
 and points of totally constructive interference called
.

\item \hfill \textit{1 pts}

Using only words from the list below, place the correct or best word in each blank in the following sentences concerning waves on a vibrating guitar string. Words may be used more than once.

Doubling the tension of a guitar string
 the speed of the wave on the string and
 the frequency of vibration. Using a string of four times the mass
 the speed of the wave by a factor of
 and
 the frequency by the same factor.



Word List: \textbf{doubles; increases; decreases; one-half; one-third; one-fourth; triples; quadruples}

\end{enumerate}

\end{document}
